\documentclass[UTF8,11pt,a4paper]{ctexart}

\usepackage[a4paper,margin=2.2cm]{geometry}
\usepackage{booktabs,longtable,tabularx,array}
\usepackage{graphicx}
\usepackage{xcolor}
\usepackage{enumitem}
\usepackage{microtype}
\usepackage[hidelinks]{hyperref}
\usepackage[numbers,sort&compress]{natbib}

\setlength{\parindent}{2em}
\setlength{\parskip}{0.25em}
\setlist{nosep,leftmargin=2em}
\renewcommand{\arraystretch}{1.18}
\newcolumntype{Y}{>{\raggedright\arraybackslash}X}
\definecolor{boundaryblue}{RGB}{235,244,252}

\title{视觉语言相册检索项目个人技术总结}
\author{张添翼\\学号：U202315231}
\date{2026年8月}

\begin{document}

\pagestyle{plain}
\maketitle

\begin{abstract}
本报告总结我在视觉语言相册检索课程项目中的个人工作。项目的开发起点是队友已经提交的
M3--M5 多路索引、检索服务与 RRF 融合基线；我没有把这些已有代码记为个人独立成果，
而是在其上继续扩展无标注开发流程、M5 融合对照、M6 视觉重排、M7 故事与缺图补全、
M4 查询理解后端以及 A7 图像编码器适配。全部本地联调在 RTX 4060 Laptop 8GB 环境下
开展，重点考察功能契约、异常降级、显存与时延。队友生产的 Qwen3.5 canonical 标注
已经完成只读验收，我在该交付上完成 adapter 迁移、Train 1993/Val 369 全量三路索引
联调、M5→M6 接口和三个 M7 故事；另整理了 100 条已审核查询及五类均衡的 50 查询
候选池。候选级相关性判断（relevance judgments）尚未冻结，因此本文只报告工程与
单查询指标链路 smoke，不据来源图正例声明正式 Recall、MRR、mAP 或 nDCG 提升。

\noindent\textbf{关键词：}视觉语言模型；多模态检索；重排序；视觉故事；本地部署
\end{abstract}

\noindent\fcolorbox{blue!55!black}{boundaryblue}{%
  \parbox{0.94\linewidth}{\textbf{贡献归属声明：}
  Git 提交 \texttt{222f232} 中原始 M3--M5 多模态检索流水线归属队友。
  本报告所称“本人工作”仅指从 \texttt{4457459} 起对该基线所做的扩展、补全、
  联调、验证与修复。}}

\section{项目背景与个人任务边界}

项目目标是让用户用中文自然语言检索个人相册，并在召回结果上完成引用问答和图文故事。
技术计划将主链路划分为：M3 建立图像、结构化文本和 BM25 索引；M4 将查询拆成语义文本
与硬条件；M5 融合三路召回；M6 用视觉语言模型重排 Top-20；M7 输出带引用回答、排序
故事及带“AI 生成”标识的补图。图像塔采用 Chinese-CLIP 或 jina-clip-v2，二者均把
文本与图片映射到同一向量空间\citep{chineseclip2022,jinaclipv2}；向量索引由 FAISS
支持\citep{faiss2017}。

系统中的模型承担不同职责。Chinese-CLIP 是默认的快速广召回编码器：它预先编码图片，
查询时只需编码文本并做向量近邻搜索。Qwen3-VL-2B-Instruct 则是较慢但能生成文本的
图像理解模型，用于候选重排、带引用问答、故事和文生图提示词\citep{qwen3vl}。
因此 CLIP 解决“从大量图片中快速找候选”，Qwen3-VL 解决“对少量候选进行深入理解”。

\begin{table}[htbp]
\centering
\small
\caption{模块、开发起点与个人任务边界}
\begin{tabularx}{\textwidth}{>{\bfseries}p{1.0cm}YYY}
\toprule
模块 & 队友已有/共同起点 & 本人扩展 & 当前结论边界 \\
\midrule
M3 & 队友已有三路索引框架 & 补全迁移、image-only 与一致性验证 & Train 1993/Val 369 三路索引已验收 \\
M4 & 检索服务接口 & 补全规则、本地 Qwen、兼容 API 三后端 & 三查询只验证接口契约 \\
M5 & 队友已有 RRF 基线 & 加入 weighted、Top-20 接口和评测池 & 候选级 qrels 未冻结 \\
M6 & 候选重排接口 & 补全 Top-20 listwise、校验和质量接口 & 单查询 smoke 不是最终结论 \\
M7 & 检索结果与 UI 起点 & 扩展自动故事、缺图补全和 AI 标识 & 三故事已跑，质量待人工评分 \\
A7 & Chinese-CLIP 默认图像塔 & 适配 jina-clip-v2 并修复 NaN & 64 图只做资源对比 \\
\bottomrule
\end{tabularx}
\end{table}

\section{队友代码基线分析}

队友在提交 \texttt{222f232} 中已经形成 M3--M5 基线：图片向量、结构化描述向量与
关键词 BM25 三路索引分别返回候选，搜索服务再用 RRF 按名次合并结果。RRF 不要求三路
原始分数处在同一尺度，适合异构召回\citep{rrf2009}。该结构提供了清楚的编码器、索引、
融合与服务边界，使后续功能可以逐层替换并保留默认行为。我的工作以此为起点，重点补足
无标注阶段的可运行入口、替代算法对照、生成式模型联调、异常防护和可复现实验，而不是
重新声明已有流水线的所有权。

\section{个人技术工作}

\subsection{无标注开发模式与任意目录流水线}

\textbf{问题：}队友仍在生成 M1 结构化标注时，完整文本索引与正式评测不能启动；若所有
开发都等待交付，接口和显存问题会集中到最后阶段。\textbf{本人工作：}我补全 mock、
image-only、full 三种运行模式，并增加任意图片目录的一键清单生成、配置复制、图像建库
和服务入口。没有标注时只运行图像检索，不伪造结构化字段；收到标注后可切换 full 模式。
\textbf{验证：}隐藏目录与 image-only 流程经过自动化测试，阶段产物均放在 Git 忽略目录。
\textbf{限制：}mock 只验证接口，image-only 只验证图像分支，两者均不能替代完整 M3 或
检索质量实验。
收到队友 Qwen3.5 canonical v1.3 后，我没有修改其原始标注，而是选择性迁移 adapter
与导入 CLI，并用 manifest 的 image ID、split、路径和 SHA-256 做只读对齐。当前已导入
Train 1993、Val 369 条记录，并在两个 split 上构建 image/text/BM25 三路索引；这项
工作属于基线迁移、联调和验收，不把队友的标注生产或原始 M3--M5 实现记为本人工作。

\subsection{M5：归一化加权融合对照}

\textbf{问题：}原有 RRF 能规避分数量纲差异，但技术计划同时要求与归一化加权和对照。
\textbf{本人工作：}我在不改变默认 RRF 的前提下，增加分支内 min--max 归一化和可配置
权重；分支不可用时只在有效分支间重新归一化权重，并保留每路原始分数、名次和降级信息。
\textbf{验证：}在 20 图、17 条有效结构化标注和 12 条查询上，两种融合均完成 Top-8
检索，候选重合率和位次变化可被脚本复现。\textbf{限制：}该对照没有人工 qrels，
92.71\% 重合率只描述两个算法输出的相似程度，不代表准确率。

\subsection{M6：Top-20 Listwise 视觉重排}

\textbf{问题：}逐图 pointwise 调用稳定，但 Top-20 需要 20 次模型推理；listwise 一次
处理候选更快，却可能输出未知、重复或遗漏 ID。多模态大模型作为重排器已有相关探索
\citep{ragvl}。\textbf{本人工作：}我把候选缩放为 5 列 contact sheet，保留编号和
image ID，并实现输出白名单、编号映射、去重、遗漏项按原序补尾以及硬失败/部分降级分开
计数。\textbf{验证：}阶段性 3 条 Top-20 查询在 8GB 显存上完成；后续 M5→M6 正式接口
接收 12 查询、240 个唯一候选，listwise 总耗时 110.968 s、峰值约 4.10 GiB，硬失败
为 0，但 8/12 出现重复或遗漏 ID 后的部分降级。\textbf{限制：}另有 q001 固定候选
单查询 smoke 验证了 MRR/nDCG 计算链路；它只有来源图正例，不能说明 listwise 排序
质量优于 pointwise，也不能替代 50 查询候选级 qrels。

\subsection{M7：自动故事与缺图补全}

\textbf{问题：}技术计划要求把 3--8 张图片组织为故事，并对叙事缺口生成插图；生成内容
必须与实拍图明确区分。视觉故事任务可参考 VIST 的图像序列叙事设定\citep{vist2016}。
\textbf{本人工作：}我扩展了按时间桶和场景相似度排序、跨时段/低相似跳变检测、Qwen
故事与提示词生成、Stable Diffusion 1.5 FP16 补图，以及
\texttt{source=generated}、\texttt{ai\_generated=true} 双重标识。进入扩散模型前卸载
Qwen，生成失败时保留原故事并显示原因。交互页面基于 Gradio\citep{gradio}。
\textbf{验证：}从 12 查询 M6 输出中选择三组、每组 5 张图片完成真实故事；其中 q03
检测到一个转场缺口并生成 512$\times$512 PNG，记录同时具有
\texttt{source=generated} 与 \texttt{ai\_generated=true}。\textbf{限制：}SD 1.5
是适合本机联调的阶段模型，并非计划首选的
SDXL/FLUX；现有结果不代表故事连贯性或风格一致性得分。

\subsection{M4：三后端查询理解}

\textbf{问题：}否定、数量和 OCR 查询不能只依赖向量相似度；同时，本地模型或免费 API
不可用时不能阻断检索。\textbf{本人工作：}我补全 rules、local Qwen 和
OpenAI-compatible 三后端。规则层先保护场景、时间、数量、OCR 和否定等硬条件；LLM 仅
补充软字段与 \texttt{semantic\_text}，失败则记录原因并回退规则。兼容 API 客户端实现
环境变量取 Key、超时、重试和响应解析。\textbf{验证：}本地 Qwen 的否定、数量、OCR
三类查询均完成结构化输出；无 Key 时 API 后端明确记录错误并安全回退。
\textbf{限制：}尚无外部服务 Key，故只完成协议测试而未声称真实免费 API 调用成功；
3/3 也不是查询理解准确率。

\subsection{A7：jina-clip-v2 适配与稳定性修复}

\textbf{问题：}A7 要求对比两个图像编码器，但 jina-clip-v2 初次加载后图像向量全为
NaN，不能直接写入 FAISS。\textbf{本人工作：}我增加编码器配置切换、Matryoshka 维度
截断与索引元数据，并用逐层钩子定位 NaN 首次出现于视觉塔第 0 层 RoPE。其非持久
\texttt{freqs\_cos/freqs\_sin} 未从 checkpoint 恢复，我按模型配置确定性重建 buffer，
同时在输出端加入 NaN/Inf 硬校验。\textbf{验证：}修复后的 8 图冒烟和 64 图资源测试均
产生有限向量与非空 Top-5；当前完整仓库回归为 248 passed。
\textbf{限制：}两模型分数不在同一标尺；资源差异不能推出检索质量差异。

\subsection{正式评测准备与防循环门禁}

\textbf{问题：}只把每条查询的来源图片标成 2，会使所有系统都容易获得正例，却无法
评价其余候选是相关、部分相关还是不相关；直接用检索分数生成标签又会形成循环论证。
\textbf{本人工作：}我合并并复核两批查询，形成 100 条来源正例集合，类别分布为
simple/compositional/negative/count/OCR = 28/26/14/14/18；再从五类各选 10 条，汇集
CLIP-only、text-only、BM25-only、三路 RRF 和 weighted 的 Top-5，形成平均 10.8 个
候选的 50 查询池。我补全候选级 0/1/2 审核页、零等级保留、qrels 冻结门禁、A5 五变体
与 A6 固定候选质量接口。\textbf{当前边界：}100 条查询已经审核，但 50 查询候选池中
非来源候选仍需逐图判断；正式结果文件不存在时，报告不生成 A5/A6 最终表。

\section{本地实验与结果}

所有数据均来自 RTX 4060 Laptop 8GB 本地阶段实验。表中“峰值显存”指 PyTorch 的
CUDA 最大分配量，而非整机 \texttt{nvidia-smi} 占用。

\begin{table}[htbp]
\centering
\small
\caption{M6 在 3 条查询、每条 Top-20 上的热态资源对照}
\begin{tabular}{lrr}
\toprule
指标 & Pointwise & Listwise \\
\midrule
模型调用数 & 60 & 3 \\
平均每查询延迟（s） & 44.377 & 9.143 \\
峰值 CUDA 显存（GiB） & 4.039 & 4.092 \\
硬失败率 & 0\% & 0\% \\
部分降级 & 不适用 & 1/3 \\
\bottomrule
\end{tabular}
\end{table}

M6 数据说明 contact-sheet 方案能在 8GB 上运行，并减少本地调用次数与这组测试的
耗时；它没有证明排序更相关。M4 中本地 Qwen 完成 3/3 条接口联调：首条冷启动
10.085 s，后两条热调用均值约 2.697 s。M5 小样共 12 查询，RRF 与 weighted 的平均
Top-8 重合率为 92.71\%，共同结果平均位次变化 0.446；单次平均延迟分别为 10.58 ms
与 6.62 ms。样本和计时轮数有限，不能据此断言 weighted 总是更快。

在正式 A6 接口上，我对 q001 的同一 Top-20 候选做了一次 source-positive smoke。
Pointwise 20 次调用耗时 37.800 s、峰值 4.04 GiB；listwise 1 次调用耗时 9.158 s、
峰值 4.09 GiB，无失败或降级，实测加速 4.128 倍。表中质量值只由来源图
\texttt{val-2007} 一个正例计算，用于验证三种顺序的指标接线，不能推广到正式集合。

\begin{table}[htbp]
\centering
\small
\caption{q001 固定 Top-20 的 source-positive 指标链路 smoke（非最终结果）}
\begin{tabular}{lrrr}
\toprule
指标 & M5 baseline & Pointwise & Listwise \\
\midrule
MRR & 0.500 & 1.000 & 0.500 \\
nDCG@10 & 0.631 & 1.000 & 0.631 \\
模型调用数 & 0 & 20 & 1 \\
推理耗时（s） & 0 & 37.800 & 9.158 \\
\bottomrule
\end{tabular}
\end{table}

\begin{table}[htbp]
\centering
\small
\caption{A7 在相同 64 张图片、512 维设置下的资源对照}
\begin{tabular}{lrr}
\toprule
指标 & Chinese-CLIP & jina-clip-v2 \\
\midrule
冷启动建库（s） & 6.060 & 10.510 \\
峰值 CUDA 显存（GiB） & 0.394 & 2.577 \\
平均热查询（ms） & 2.251 & 47.701 \\
输出维度 & 512 & 512 \\
\bottomrule
\end{tabular}
\end{table}

在本次 batch 设置下，Chinese-CLIP 的建库、显存与查询资源开销较低，因此继续作为
默认快速召回编码器；jina-clip-v2 保留为多语言和可截断表示的 A7 对照。此结论仅针对
资源，正式选择仍需在同一人工 qrels 上比较 Recall@K、MRR、mAP 与 nDCG@10。M7 的
Qwen3-VL-2B 故事链及 SD 1.5 FP16 缺图生成也在本机串行运行，证明链路可执行，未给出
故事质量分数。

\section{关键问题、修复与反思}

\begin{enumerate}
  \item \textbf{结构化输出不等于结构可靠。}Qwen 曾把数组输出成字符串，还把 prompt
  示例中的否定词写入查询。修复不是继续堆示例，而是在入口做类型归一化、收紧 schema，
  并让硬过滤只信任确定性规则。由此我认识到 LLM 输出必须被当作不可信外部输入。
  \item \textbf{部分结果也要可审计。}Listwise 曾重复 2 个 ID、遗漏 3 个 ID。系统去重
  并补尾可以保持页面可用，但必须标记 \texttt{degraded=true}，否则会掩盖模型失约。
  \item \textbf{数值正常性应在索引前阻断。}Jina 的 NaN 并非模型权重损坏，而是低内存
  加载下非持久 EVA RoPE buffer 未初始化。逐层定位、确定性重建和输出端有限值校验共同
  构成修复；只在最终检索结果上观察会更难定位。
  \item \textbf{8GB 的关键是生命周期而非单个峰值。}我采用小 batch、缩略 contact
  sheet、FP16、attention/VAE slicing，并按“检索编码器$\rightarrow$Qwen$\rightarrow$
  Stable Diffusion”串行卸载加载，避免多个模型同时常驻。该策略以时延换取本机可运行性。
\end{enumerate}

这些问题也说明，自动化测试通过与正式实验完成是两件事：当前 248 项测试通过，
证明接口和回归状态，却不能替代候选级人工相关性、数据覆盖和统计显著性。

\section{标注到达后的完成项与剩余门禁}

队友标注到达后，我按只读边界完成了以下联调：

\begin{enumerate}
  \item 验收 Qwen3.5 canonical 2362 条输入；7 个 Train ID 缺失被如实保留，不由程序
  合成，也不回退到其他模型。
  \item 导入 Train 1993、Val 369 条统一标注，并构建两个 split 的 image/text/BM25
  三路索引；manifest 保存记录数、维度、模型指纹、配置与标注哈希。
  \item 冻结 M5→M6 v1.0 接口，验证 12 条查询、每条 20 个唯一候选；完成 listwise
  资源测试并显式披露 8/12 部分降级。
  \item 完成三个 M7 故事与一个真实缺图生成，验证故事顺序、失败保留和 AI provenance。
  \item 审核 100 条五类查询，生成 50 查询多系统候选池，补全 A5/A6 正式结果代码。
\end{enumerate}

当前剩余门禁只有候选级人工判断及其下游结果冻结：必须把候选池中每张非来源图片明确
标为 0/1/2，要求 50 个 \texttt{graded\_query\_ids} 完整、保留零等级，再运行 A5 五变体
和 A6 baseline/pointwise/listwise。随后才能把正式 MRR、mAP、nDCG、失败率和分类分析
写入报告。第二审核者尚未实际签字时，也必须保留单审核者限制。

\section{个人总结与贡献审计}

我的主要贡献不是重写队友已有的 M3--M5，而是把基线扩展成在无标注期可开发、在 8GB
本机可联调、出现生成式模型异常时可降级、标注到达后可切换到正式实验的系统。我完成了
选择性迁移、接口对照、M6/M7 生成式链路、M4 多后端、A7 第二编码器、正式评测工具、
部署与 Demo 文档，以及相应测试和故障定位。当前最重要的边界是：全量链路已经具备，
正式 A5/A6 质量结论仍须完成 50 查询候选级人工判断。

\begin{table}[htbp]
\centering
\small
\caption{个人贡献审计}
\begin{tabularx}{\textwidth}{YYY}
\toprule
队友已有 & 本人新增 & 共同依赖/后续 \\
\midrule
M3--M5 三路索引、搜索服务、RRF 基线 & Qwen3.5 选择性迁移、全量索引验收、weighted 对照 & M1 标注生产仍归属队友 \\
候选检索和服务边界 & M6 Top-20 listwise、接口校验、8GB 与指标接线 & 50 查询候选级 qrels \\
项目 UI 与输出接口起点 & M7 三故事、SD 补图、AI 标识及页面联调 & 故事人工评价和最终录屏 \\
Chinese-CLIP 图像塔起点 & M4 三后端；A7 Jina；正式评测与部署 & A5/A6 冻结表与团队报告 \\
\bottomrule
\end{tabularx}
\end{table}

\bibliographystyle{plainnat}
\bibliography{references}

\appendix
\section{提交与证据映射}

\small
\begin{longtable}{p{3.4cm}p{9.6cm}}
\toprule
提交 & 对应证据 \\
\midrule
\endfirsthead
\toprule
提交 & 对应证据 \\
\midrule
\endhead
\texttt{222f232} & 队友提交的 M3--M5 多模态检索基线；仅作为本人开发起点。 \\
\texttt{4457459} & 补全无标注开发模式、image-only 和评测门禁。 \\
\texttt{9c5d54e} & 扩展路由、人工评测工具、M6 基准与 M7 UI 桥接。 \\
\texttt{33551fe} & 补全任意图片目录的一键流水线与录屏指南。 \\
\texttt{6550040} & 扩展 M5 归一化加权融合及 RRF 工程对照。 \\
\texttt{c4d329c} & 补全 M6 Top-20 listwise 重排、降级逻辑和 8GB 基准。 \\
\texttt{0c54bf8} & 扩展 M7 自动故事、缺口检测、补图和 AI 生成标识。 \\
\texttt{223f3e0} & 补全 M4 三后端、回退与结构化输出验证。 \\
\texttt{a8f3a05} & 扩展 A7 jina-clip-v2、修复 RoPE/NaN 并完成资源对比。 \\
\texttt{bfc2c1c} & 迁移队友 Qwen3.5 adapter 与只读导入 CLI，不声明标注生产归属。 \\
\texttt{fc4e7c6} & 补全 Train/Val 全量三路索引一致性和 manifest 验证。 \\
\texttt{e7f40da--557216b} & 规定并验证 M5→M6 接口，补全离线 listwise 与资源指标。 \\
\texttt{94250c1,82478ca,}\newline\texttt{9451b78} & 补全 canonical M7 桥接、三个故事和正式运行记录。 \\
\texttt{c64a958} & 修复正式 qrels 的零等级保留。 \\
\texttt{6deb435} & 整理两批已审核查询为正式 100 条集合。 \\
\texttt{7401f86} & 构建五类均衡、五个 A5 变体的 50 查询候选池。 \\
\texttt{d5422a1} & 补全候选级 0/1/2 图片审核页。 \\
\texttt{ded5d2a} & 补全 qrels 冻结与完整性门禁。 \\
\texttt{a784a5e} & 补全 A5 五变体正式输出与来源哈希。 \\
\texttt{6b2401a} & 补全 A6 固定候选质量接口并完成 q001 smoke。 \\
\texttt{c9873fb,260fb4c} & 补全可复现 GPU 部署和最终 Demo 手册。 \\
\bottomrule
\end{longtable}

\end{document}
